% 原书第 77—79 页；第 5.2—5.3 节。
\section{协变量偏移}
\label{sec:ch05-covariate-shift}

如第 2.1 节所述，$P_{\mathrm{cause}}$ 与 $P_{\mathrm{effect}\mid\mathrm{cause}}$ 之间的独立性，即原理 2.1，可以从两种不同的角度理解。在上面的第 5.1 节中，我们说明，对于固定的联合分布，这两个对象不包含关于彼此的信息，见图 2.2 中间的方框。另一种理解是：假设联合分布 $P_{\mathrm{cause},\mathrm{effect}}$ 在不同数据集之间发生变化，那么 $P_{\mathrm{cause}}$ 的变化，并不会告诉我们 $P_{\mathrm{effect}\mid\mathrm{cause}}$ 如何变化，这对应于图 2.2 左侧的方框。因此，在根据 $X$ 预测 $Y$ 的情景中，知道 $X$ 是原因、$Y$ 是结果，会产生重要影响。

假设我们利用一个数据集中的样本，已经学到了 $X$ 与 $Y$ 之间的统计关系，现在要用这些知识，在第二个数据集中根据 $X$ 预测 $Y$。进一步假设，我们发现第二个数据集中的 $x$ 值服从分布 $P'_X$，与第一个数据集的分布 $P_X$ 不同。应当如何利用这一信息？根据机制独立性，$P'_X$ 不同于 $P_X$，并不能告诉我们，$P_{Y\mid X}$ 在两个数据集之间是否也发生了变化。因此，条件分布 $P_{Y\mid X}$ 仍可能适用于第二个数据集。其次，即使条件分布确实发生了变化，即 $P'_{Y\mid X}\neq P_{Y\mid X}$，继续用 $P_{Y\mid X}$ 预测，仍然是自然的做法。毕竟，机制独立原理指出，边缘分布由 $P_X$ 变成 $P'_X$，并不能告诉我们条件分布是如何变化的。因此，在没有更好选择的情况下，我们仍使用 $P_{Y\mid X}$。尽管 $P_X$ 已发生变化，却继续使用相同的条件分布 $P_{Y\mid X}$，通常称为\textbf{协变量偏移}。目前，机器学习领域已对这一假设作了广泛研究［Sugiyama and Kawanabe, 2012］。Schölkopf et al.［2012］指出，这种做法只有在沿因果方向预测时，即 $X$ 是原因、$Y$ 是结果时，才具有依据。

\Needspace{6\baselineskip}
为了进一步说明这一点，考虑一个沿反因果方向预测的简化例子，其中 $X$ 是结果。设 $Y$ 是一个二元变量，以加性的方式影响实值变量 $X$：
\begin{equation}
 X=Y+N_X,
 \tag{5.4}
\end{equation}
其中，假定 $N_X$ 是独立于 $Y$ 的高斯噪声。图 5.4 左图展示了相应的概率密度 $p_X$。

\begin{figure}[htbp]
\centering
\begin{tikzpicture}[x=1cm,y=1cm]
\begin{scope}
\draw[thin] (0,2.4)--(0,0)--(6,0);
% Original printed PDF vector path.
\path[draw=plotgray,line width=.5pt,fill=gray!55] (0.50002,0.20001) .. controls (0.50002,0.20001) and (0.57375,0.20002) ..(0.60206,0.20003) .. controls (0.63038,0.20004) and (0.67579,0.20006) ..(0.70410,0.20011) .. controls (0.73242,0.20016) and (0.77783,0.20023) ..(0.80614,0.20038) .. controls (0.83446,0.20053) and (0.87987,0.20077) ..(0.90818,0.20119) .. controls (0.93650,0.20161) and (0.98191,0.20233) ..(1.01022,0.20343) .. controls (1.03854,0.20453) and (1.08395,0.20650) ..(1.11226,0.20914) .. controls (1.14058,0.21177) and (1.18599,0.21665) ..(1.21430,0.22242) .. controls (1.24262,0.22819) and (1.28803,0.23916) ..(1.31634,0.25072) .. controls (1.34466,0.26228) and (1.39007,0.28460) ..(1.41838,0.30574) .. controls (1.44670,0.32689) and (1.49211,0.36789) ..(1.52042,0.40314) .. controls (1.54874,0.43840) and (1.59415,0.50656) ..(1.62246,0.55985) .. controls (1.65078,0.61314) and (1.69619,0.71460) ..(1.72450,0.78722) .. controls (1.75282,0.85985) and (1.79823,0.99485) ..(1.82654,1.08328) .. controls (1.85486,1.17171) and (1.90027,1.33005) ..(1.92858,1.42456) .. controls (1.95690,1.51907) and (2.00231,1.67880) ..(2.03062,1.76446) .. controls (2.05894,1.85013) and (2.10435,1.98167) ..(2.13266,2.04198) .. controls (2.16098,2.10230) and (2.20639,2.17726) ..(2.23470,2.19918) .. controls (2.26302,2.22110) and (2.30842,2.22128) ..(2.33674,2.19999) .. controls (2.36506,2.17870) and (2.41046,2.10481) ..(2.43878,2.04572) .. controls (2.46710,1.98664) and (2.51250,1.85714) ..(2.54082,1.77415) .. controls (2.56914,1.69115) and (2.61454,1.53630) ..(2.64286,1.44756) .. controls (2.67118,1.35882) and (2.71658,1.21147) ..(2.74490,1.13457) .. controls (2.77322,1.05767) and (2.81862,0.94488) ..(2.84694,0.89332) .. controls (2.87526,0.84175) and (2.92066,0.78104) ..(2.94898,0.76295) .. controls (2.97730,0.74487) and (3.02270,0.74495) ..(3.05102,0.76298) .. controls (3.07933,0.78100) and (3.12474,0.84143) ..(3.15306,0.89289) .. controls (3.18137,0.94435) and (3.22678,1.05703) ..(3.25510,1.13386) .. controls (3.28341,1.21069) and (3.32882,1.35787) ..(3.35714,1.44660) .. controls (3.38545,1.53533) and (3.43086,1.69029) ..(3.45918,1.77336) .. controls (3.48749,1.85643) and (3.53290,1.98615) ..(3.56122,2.04531) .. controls (3.58953,2.10447) and (3.63494,2.17840) ..(3.66326,2.19977) .. controls (3.69157,2.22114) and (3.73698,2.22117) ..(3.76530,2.19936) .. controls (3.79361,2.17756) and (3.83902,2.10282) ..(3.86734,2.04259) .. controls (3.89565,1.98236) and (3.94106,1.85089) ..(3.96938,1.76525) .. controls (3.99769,1.67961) and (4.04310,1.51986) ..(4.07142,1.42535) .. controls (4.09973,1.33083) and (4.14514,1.17253) ..(4.17346,1.08408) .. controls (4.20177,0.99563) and (4.24718,0.86053) ..(4.27550,0.78785) .. controls (4.30381,0.71517) and (4.34922,0.61360) ..(4.37754,0.56026) .. controls (4.40585,0.50691) and (4.45126,0.43863) ..(4.47958,0.40335) .. controls (4.50789,0.36806) and (4.55330,0.32709) ..(4.58162,0.30592) .. controls (4.60993,0.28476) and (4.65534,0.26240) ..(4.68366,0.25082) .. controls (4.71197,0.23924) and (4.75738,0.22825) ..(4.78570,0.22247) .. controls (4.81401,0.21669) and (4.85942,0.21180) ..(4.88774,0.20916) .. controls (4.91605,0.20652) and (4.96146,0.20454) ..(4.98978,0.20344) .. controls (5.01809,0.20233) and (5.06350,0.20162) ..(5.09182,0.20119) .. controls (5.12013,0.20077) and (5.16554,0.20053) ..(5.19386,0.20038) .. controls (5.22217,0.20023) and (5.26758,0.20016) ..(5.29589,0.20011) .. controls (5.32421,0.20006) and (5.36962,0.20005) ..(5.39793,0.20003) .. controls (5.42625,0.20002) and (5.49998,0.20001) ..(5.49998,0.20001);

\node[below] at (3,-.05){$x$};
\node[rotate=90] at (-.42,1.2){$p(x)$};
\end{scope}
\begin{scope}[xshift=7.4cm]
\draw[thin] (0,2.4)--(0,0)--(6,0);
% Original printed PDF vector path.
\path[draw=plotgray,line width=.5pt,fill=gray!55] (0.50002,0.20002) .. controls (0.50002,0.20002) and (0.57375,0.20003) ..(0.60206,0.20004) .. controls (0.63038,0.20006) and (0.67579,0.20007) ..(0.70410,0.20012) .. controls (0.73242,0.20017) and (0.77783,0.20024) ..(0.80614,0.20039) .. controls (0.83446,0.20054) and (0.87987,0.20077) ..(0.90818,0.20120) .. controls (0.93650,0.20162) and (0.98191,0.20234) ..(1.01022,0.20344) .. controls (1.03854,0.20455) and (1.08395,0.20651) ..(1.11226,0.20915) .. controls (1.14058,0.21179) and (1.18599,0.21667) ..(1.21430,0.22244) .. controls (1.24262,0.22821) and (1.28803,0.23919) ..(1.31634,0.25075) .. controls (1.34466,0.26232) and (1.39007,0.28464) ..(1.41838,0.30580) .. controls (1.44670,0.32696) and (1.49211,0.36797) ..(1.52042,0.40325) .. controls (1.54874,0.43852) and (1.59415,0.50671) ..(1.62246,0.56003) .. controls (1.65078,0.61334) and (1.69619,0.71485) ..(1.72450,0.78750) .. controls (1.75282,0.86016) and (1.79823,0.99522) ..(1.82654,1.08369) .. controls (1.85486,1.17216) and (1.90026,1.33057) ..(1.92858,1.42513) .. controls (1.95690,1.51968) and (2.00231,1.67948) ..(2.03062,1.76517) .. controls (2.05894,1.85087) and (2.10434,1.98245) ..(2.13266,2.04276) .. controls (2.16098,2.10307) and (2.20639,2.17801) ..(2.23470,2.19983) .. controls (2.26302,2.22165) and (2.30842,2.22165) ..(2.33674,2.20003) .. controls (2.36505,2.17841) and (2.41046,2.10395) ..(2.43878,2.04402) .. controls (2.46710,1.98408) and (2.51250,1.85304) ..(2.54082,1.76804) .. controls (2.56914,1.68303) and (2.61454,1.52445) ..(2.64286,1.43134) .. controls (2.67117,1.33823) and (2.71658,1.18258) ..(2.74490,1.09698) .. controls (2.77322,1.01137) and (2.81862,0.88177) ..(2.84694,0.81434) .. controls (2.87525,0.74691) and (2.92066,0.65549) ..(2.94898,0.61096) .. controls (2.97730,0.56643) and (3.02270,0.51535) ..(3.05102,0.49340) .. controls (3.07933,0.47145) and (3.12474,0.45577) ..(3.15306,0.45276) .. controls (3.18137,0.44975) and (3.22678,0.46116) ..(3.25510,0.47170) .. controls (3.28341,0.48223) and (3.32882,0.51080) ..(3.35714,0.52866) .. controls (3.38545,0.54652) and (3.43086,0.58161) ..(3.45918,0.60040) .. controls (3.48749,0.61920) and (3.53290,0.65017) ..(3.56122,0.66414) .. controls (3.58953,0.67811) and (3.63494,0.69606) ..(3.66326,0.70108) .. controls (3.69157,0.70611) and (3.73698,0.70594) ..(3.76530,0.70038) .. controls (3.79361,0.69481) and (3.83902,0.67606) ..(3.86734,0.66096) .. controls (3.89565,0.64586) and (3.94106,0.61296) ..(3.96938,0.59153) .. controls (3.99769,0.57010) and (4.04310,0.53014) ..(4.07142,0.50650) .. controls (4.09973,0.48286) and (4.14514,0.44326) ..(4.17346,0.42114) .. controls (4.20177,0.39901) and (4.24718,0.36522) ..(4.27550,0.34704) .. controls (4.30381,0.32886) and (4.34922,0.30346) ..(4.37754,0.29012) .. controls (4.40585,0.27677) and (4.45126,0.25970) ..(4.47958,0.25087) .. controls (4.50789,0.24205) and (4.55330,0.23180) ..(4.58162,0.22650) .. controls (4.60993,0.22121) and (4.65534,0.21562) ..(4.68366,0.21272) .. controls (4.71197,0.20983) and (4.75738,0.20708) ..(4.78569,0.20563) .. controls (4.81401,0.20419) and (4.85942,0.20296) ..(4.88774,0.20230) .. controls (4.91605,0.20164) and (4.96146,0.20115) ..(4.98978,0.20087) .. controls (5.01809,0.20059) and (5.06350,0.20041) ..(5.09181,0.20031) .. controls (5.12013,0.20020) and (5.16554,0.20014) ..(5.19386,0.20011) .. controls (5.22217,0.20007) and (5.26758,0.20005) ..(5.29589,0.20004) .. controls (5.32421,0.20003) and (5.36962,0.20002) ..(5.39793,0.20002) .. controls (5.42625,0.20001) and (5.49997,0.20001) ..(5.49997,0.20001);

\node[below] at (3,-.05){$x$};
\node[rotate=90] at (-.42,1.2){$p'(x)$};
\end{scope}
\end{tikzpicture}
\caption{$P_X$ 变为 $P'_X$ 的一种情形：这种变化提示，$P_Y$ 已改变，而 $P_{X\mid Y}$ 保持不变。当已知二元变量 $Y$ 是 $X$ 的原因时，观测到 $P_X$ 为两个高斯分布的混合，会使我们有理由认为，两个峰分别对应不同的标签 $y=0,1$。于是，$Y$ 对 $X$ 的影响只是平移高斯分布的均值，这就构成加性噪声模型，见第 4.1.4 节。这无疑为联合分布提供了简单解释。如果进一步观测到，混合权重在两个数据集之间发生了变化，那么这一变化就很可能是由 $P_Y$ 的变化造成的。}
\label{fig:ch05-changing-mixture}
\end{figure}


只要噪声的宽度足够小，分布 $P_X$ 就具有双峰。即使完全不了解生成模型，也可以看出，$P_X$ 是两个宽度相同的高斯分布的混合。因此，在这种情况下，仅凭 $P_X$ 就能够猜测联合分布 $P_{X,Y}$，因为一个自然的假设是：$Y$ 的影响只是平移 $X$ 的均值。在这一假设下，不需要任何 $(x,y)$ 对，就能学到 $X$ 与 $Y$ 的关系。现在假设，在第二个数据集中，我们观测到的仍是这两个高斯分布的混合，但其权重不同，见图 5.4 右图。此时，最自然的结论是：同一个式 (5.4) 仍然成立，只是 $P_Y$ 改变了，因而混合权重发生了变化。相应地，我们就不会再使用相同的 $P_{Y\mid X}$ 进行预测，而会根据 $P'_X$ 重建 $P'_{Y\mid X}$。这个例子说明，在反因果情景中，$P_X$ 与 $P_{Y\mid X}$ 的变化可能相互关联，而这种关联可能源于 $P_Y$ 发生变化、$P_{X\mid Y}$ 保持不变。换言之，正因为 $P_{\mathrm{cause}}$ 与 $P_{\mathrm{effect}\mid\mathrm{cause}}$ 独立变化，$P_{\mathrm{effect}}$ 与 $P_{\mathrm{cause}\mid\mathrm{effect}}$ 往往会以相依的方式变化。

上述例子展示了一种特定情景。要构想一般的方法，利用 $P_{\mathrm{effect}}$ 与 $P_{\mathrm{cause}\mid\mathrm{effect}}$ 相依变化这一事实，仍是一个困难问题。这可能是值得进一步探索的研究方向。我们认为，因果性能够在\textbf{领域适应（domain adaptation）}和迁移问题中发挥重要作用，亦见 Bareinboim and Pearl［2016］、Rojas-Carulla et al.［2016］、Zhang et al.［2013］和 Zhang et al.［2015］。

\begin{figure}[t]
\centering
\begin{tikzpicture}[x=1cm,y=1cm]
\begin{scope}
\draw[thin] (0,2.4)--(0,0)--(6,0);
\foreach \x/\lab in {1.57143/1,2.82143/2,4.78571/3}{
 \draw[gray!50,thin] (\x,0)--(\x,.1862);
 \node[below] at (\x,-.02){$s_{\lab}$};}
\foreach \y in {.2,1.1,2.0}{\draw[gray!50,thin] (0,\y)--(.1766,\y);}
% Original printed PDF vector path.
\path[draw=plotgray,line width=.5pt] (0.50002,0.20002) .. controls (0.50002,0.20002) and (0.57375,0.20002) ..(0.60206,0.20002) .. controls (0.63038,0.20002) and (0.67579,0.20002) ..(0.70410,0.20002) .. controls (0.73242,0.20002) and (0.77783,0.20002) ..(0.80614,0.20002) .. controls (0.83446,0.20002) and (0.87987,0.20002) ..(0.90818,0.20002) .. controls (0.93650,0.20002) and (0.98191,0.19999) ..(1.01022,0.20002) .. controls (1.03854,0.20004) and (1.08395,0.19915) ..(1.11226,0.20022) .. controls (1.14058,0.20129) and (1.18599,0.19148) ..(1.21430,0.20774) .. controls (1.24262,0.22400) and (1.28803,0.21850) ..(1.31634,0.31741) .. controls (1.34466,0.41631) and (1.39007,0.68912) ..(1.41838,0.92059) .. controls (1.44670,1.15205) and (1.49211,1.83776) ..(1.52042,1.98565) .. controls (1.54874,2.13353) and (1.59415,2.13408) ..(1.62246,1.98645) .. controls (1.65078,1.83881) and (1.69619,1.15313) ..(1.72450,0.92159) .. controls (1.75282,0.69005) and (1.79823,0.41673) ..(1.82654,0.31769) .. controls (1.85486,0.21865) and (1.90027,0.22406) ..(1.92858,0.20776) .. controls (1.95690,0.19147) and (2.00231,0.20130) ..(2.03062,0.20022) .. controls (2.05894,0.19915) and (2.10434,0.20004) ..(2.13266,0.20002) .. controls (2.16098,0.19999) and (2.20639,0.20001) ..(2.23470,0.20002) .. controls (2.26302,0.20002) and (2.30842,0.19962) ..(2.33674,0.20009) .. controls (2.36506,0.20056) and (2.41046,0.19445) ..(2.43878,0.20340) .. controls (2.46710,0.21236) and (2.51250,0.19603) ..(2.54082,0.26464) .. controls (2.56914,0.33326) and (2.61454,0.49200) ..(2.64286,0.69792) .. controls (2.67118,0.90385) and (2.71658,1.54801) ..(2.74490,1.74880) .. controls (2.77322,1.94960) and (2.81862,2.22316) ..(2.84694,2.14511) .. controls (2.87525,2.06706) and (2.92066,1.42815) ..(2.94898,1.18627) .. controls (2.97730,0.94440) and (3.02270,0.53638) ..(3.05102,0.40186) .. controls (3.07933,0.26733) and (3.12474,0.24463) ..(3.15306,0.21670) .. controls (3.18137,0.18877) and (3.22678,0.20289) ..(3.25510,0.20057) .. controls (3.28341,0.19826) and (3.32882,0.20010) ..(3.35714,0.20002) .. controls (3.38545,0.19995) and (3.43086,0.20002) ..(3.45918,0.20002) .. controls (3.48749,0.20002) and (3.53290,0.20002) ..(3.56122,0.20002) .. controls (3.58953,0.20002) and (3.63494,0.20002) ..(3.66326,0.20002) .. controls (3.69157,0.20002) and (3.73698,0.20002) ..(3.76530,0.20002) .. controls (3.79361,0.20002) and (3.83902,0.20002) ..(3.86734,0.20002) .. controls (3.89565,0.20002) and (3.94106,0.20002) ..(3.96938,0.20002) .. controls (3.99769,0.20002) and (4.04310,0.20002) ..(4.07142,0.20002) .. controls (4.09973,0.20002) and (4.14514,0.20001) ..(4.17346,0.20002) .. controls (4.20177,0.20002) and (4.24718,0.19984) ..(4.27550,0.20004) .. controls (4.30381,0.20023) and (4.34922,0.19676) ..(4.37754,0.20142) .. controls (4.40585,0.20607) and (4.45126,0.18872) ..(4.47958,0.23359) .. controls (4.50789,0.27847) and (4.55330,0.35351) ..(4.58162,0.52488) .. controls (4.60993,0.69625) and (4.65534,1.23631) ..(4.68366,1.46874) .. controls (4.71197,1.70117) and (4.75738,2.19939) ..(4.78570,2.20005) .. controls (4.81401,2.20072) and (4.85942,1.70563) ..(4.88773,1.47354) .. controls (4.91605,1.24145) and (4.96146,0.69929) ..(4.98978,0.52730) .. controls (5.01809,0.35531) and (5.06350,0.27920) ..(5.09182,0.23398) .. controls (5.12013,0.18877) and (5.16554,0.20615) ..(5.19386,0.20144) .. controls (5.22217,0.19673) and (5.26758,0.20024) ..(5.29589,0.20004) .. controls (5.32421,0.19984) and (5.36962,0.20002) ..(5.39793,0.20002) .. controls (5.42625,0.20001) and (5.49998,0.20002) ..(5.49998,0.20002);

\node[below] at (3,-.65){$x$};
\node[rotate=90] at (-.42,1.2){$p(x)$};
\end{scope}
\begin{scope}[xshift=7.4cm]
\draw[thin] (0,2.4)--(0,0)--(6,0);
\foreach \x/\lab in {1.57143/1,2.82143/2,4.78571/3}{
 \draw[gray!50,thin] (\x,0)--(\x,.1862);
 \node[below] at (\x,-.02){$s_{\lab}$};}
\foreach \y in {.2,1.1,2.0}{\draw[gray!50,thin] (0,\y)--(.1766,\y);}
% Original printed PDF vector path.
\path[draw=plotgray,line width=.5pt] (0.50002,0.20000) .. controls (0.50002,0.20000) and (0.57375,0.19999) ..(0.60206,0.20006) .. controls (0.63038,0.20014) and (0.67579,0.20008) ..(0.70410,0.20055) .. controls (0.73242,0.20101) and (0.77783,0.20116) ..(0.80614,0.20339) .. controls (0.83446,0.20562) and (0.87987,0.20809) ..(0.90818,0.21661) .. controls (0.93650,0.22512) and (0.98191,0.23915) ..(1.01022,0.26477) .. controls (1.03854,0.29038) and (1.08395,0.34105) ..(1.11226,0.40122) .. controls (1.14058,0.46140) and (1.18599,0.58978) ..(1.21430,0.69847) .. controls (1.24262,0.80715) and (1.28803,1.03871) ..(1.31634,1.18456) .. controls (1.34466,1.33042) and (1.39007,1.61653) ..(1.41838,1.74970) .. controls (1.44670,1.88286) and (1.49211,2.08957) ..(1.52042,2.14436) .. controls (1.54874,2.19914) and (1.59415,2.19925) ..(1.62246,2.14458) .. controls (1.65078,2.08990) and (1.69619,1.88339) ..(1.72450,1.75028) .. controls (1.75282,1.61717) and (1.79823,1.33101) ..(1.82654,1.18520) .. controls (1.85486,1.03940) and (1.90026,0.80787) ..(1.92858,0.69943) .. controls (1.95690,0.59098) and (2.00231,0.46234) ..(2.03062,0.40363) .. controls (2.05894,0.34491) and (2.10434,0.29567) ..(2.13266,0.27625) .. controls (2.16098,0.25684) and (2.20639,0.25232) ..(2.23470,0.26368) .. controls (2.26302,0.27504) and (2.30842,0.31062) ..(2.33674,0.35815) .. controls (2.36505,0.40567) and (2.41046,0.51046) ..(2.43878,0.60621) .. controls (2.46710,0.70195) and (2.51250,0.90853) ..(2.54082,1.04820) .. controls (2.56914,1.18788) and (2.61454,1.47022) ..(2.64286,1.61287) .. controls (2.67117,1.75552) and (2.71658,1.99677) ..(2.74490,2.07634) .. controls (2.77322,2.15591) and (2.81862,2.21416) ..(2.84694,2.18637) .. controls (2.87525,2.15859) and (2.92066,1.99527) ..(2.94898,1.87609) .. controls (2.97730,1.75691) and (3.02270,1.47610) ..(3.05102,1.32742) .. controls (3.07933,1.17873) and (3.12474,0.92508) ..(3.15306,0.80448) .. controls (3.18137,0.68388) and (3.22678,0.52988) ..(3.25510,0.45822) .. controls (3.28341,0.38656) and (3.32882,0.32052) ..(3.35714,0.28801) .. controls (3.38545,0.25550) and (3.43086,0.23539) ..(3.45918,0.22390) .. controls (3.48749,0.21240) and (3.53290,0.20836) ..(3.56122,0.20517) .. controls (3.58953,0.20198) and (3.63494,0.20159) ..(3.66326,0.20089) .. controls (3.69157,0.20019) and (3.73698,0.20023) ..(3.76530,0.20014) .. controls (3.79361,0.20004) and (3.83902,0.20004) ..(3.86734,0.20021) .. controls (3.89565,0.20039) and (3.94106,0.20036) ..(3.96938,0.20140) .. controls (3.99769,0.20244) and (4.04310,0.20323) ..(4.07142,0.20770) .. controls (4.09973,0.21218) and (4.14514,0.21847) ..(4.17346,0.23367) .. controls (4.20177,0.24887) and (4.24718,0.27676) ..(4.27550,0.31724) .. controls (4.30381,0.35773) and (4.34922,0.44181) ..(4.37754,0.52544) .. controls (4.40585,0.60906) and (4.45126,0.78890) ..(4.47958,0.91994) .. controls (4.50789,1.05097) and (4.55330,1.32204) ..(4.58162,1.46983) .. controls (4.60993,1.61762) and (4.65534,1.88383) ..(4.68366,1.98512) .. controls (4.71197,2.08642) and (4.75738,2.19967) ..(4.78569,2.19989) .. controls (4.81401,2.20011) and (4.85942,2.08766) ..(4.88774,1.98670) .. controls (4.91605,1.88574) and (4.96146,1.61994) ..(4.98978,1.47223) .. controls (5.01809,1.32452) and (5.06350,1.05332) ..(5.09181,0.92212) .. controls (5.12013,0.79092) and (5.16554,0.61048) ..(5.19386,0.52664) .. controls (5.22217,0.44279) and (5.26758,0.35843) ..(5.29589,0.31781) .. controls (5.32421,0.27719) and (5.36962,0.24913) ..(5.39793,0.23386) .. controls (5.42625,0.21860) and (5.49997,0.20775) ..(5.49997,0.20775);

\node[below] at (3,-.65){$y$};
\node[rotate=90] at (-.42,1.2){$p(y)$};
\end{scope}
\end{tikzpicture}
\caption{一个由 $X$ 导致 $Y$ 的例子。由此，$P_Y$ 与 $P_{X\mid Y}$ 包含关于彼此的信息。左：$P_X$ 是由位于 $s_1,s_2,s_3$ 的尖峰组成的混合分布。右：$P_Y$ 由 $P_X$ 与零均值高斯噪声进行卷积（convolution）得到，因此，其峰仍位于 $s_1,s_2,s_3$，但不再那么尖锐。于是，$P_{X\mid Y}$ 同样包含关于 $s_1,s_2,s_3$ 的信息，见习题 5.1。}
\label{fig:ch05-shared-peak-information}
\end{figure}

\FloatBarrier
\section{习题}
\label{sec:ch05-problems}

% 本章没有先行的定理类结果，原书习题从 5.1 开始。
\setcounter{problem}{0}
\begin{problem}[机制独立性]
设 $P_X$ 是由 $k$ 个尖锐高斯峰组成的混合分布，峰的位置为 $s_1,\ldots,s_k$，如图 5.5 左图所示。向 $X$ 加上均值为零、宽度为 $\sigma_N$ 的高斯噪声 $N$，得到 $Y$。噪声的宽度应使各个峰仍然清晰可辨，如图 5.5 右图所示。
\begin{enumerate}[label=\alph*)]
 \item 从直观上说明，为什么 $P_{X\mid Y}$ 也包含峰位置 $s_1,\ldots,s_k$ 的信息，因而 $P_{X\mid Y}$ 与 $P_Y$ 共享这些信息。
 \item $P_X$ 与 $P_Y$ 之间的转换，可以用卷积（从 $P_X$ 到 $P_Y$）和反卷积（deconvolution；从 $P_Y$ 到 $P_X$）描述。如果将 $P_{Y\mid X}$ 视为把输入 $P_X$ 转换为输出 $P_Y$ 的线性映射，那么 $P_{Y\mid X}$ 就是卷积映射。说明为什么 $P_{X\mid Y}$ 并不是人们乍看之下可能以为的反卷积映射。
\end{enumerate}
\end{problem}
\FloatBarrier
